\chapter[Sermon 23. The Second Vision Is Shown to St. John, Wherein He Sees God in His Throne with Elders, Whom He Describes Gallantly.]{The Second Vision Is Shown to St. John, Wherein He Sees God in His Throne with Elders, Whom He Describes Gallantly.}
\label{sermon:023}
\markright{Sermon 23. The Second Vision Is Shown to St. John, Wherein He ...}

After this, I looked, and behold, a door was open in heaven, and the first voice which I heard was like a trumpet speaking with me, which said, “Come up here, and I will show you things which must be fulfilled hereafter.” And immediately I was in the spirit, and behold, a throne was set in heaven, and one sat on the throne. And he who sat appeared like a jasper stone and a sardine stone; and there was a rainbow around the throne, in appearance like an emerald. And around the throne were twenty-four seats. And upon the seats were twenty-four elders sitting, clothed in white raiment, and had on their heads crowns of gold.

\index{Antichrist}The third part of this work reaches from the beginning of chapter 4 to the beginning of chapter 12. It contains a notable vision, most wholesome, and of much fruit. The first vision, which we heard explained in chapter 3, exhibits a figure of Christ and of his church, and how the Lord reigns in them, and how the church behaves or ought to conduct herself. In the second vision, John declares how, by a most just and most holy government, God governs all things through Christ, and what happens to the church in the world and of the world. In these are rehearsed the most sorrowful destinies of the church, calamities, plagues, and destructions, famines, persecutions, revolts, heresies, conflicts, and other evils of the same sort. \&c. Who also, and what, and how just God is in all his judgments, is described here: That he is the Author of all; that God through the most astute and excellent government of Christ rules all things; that the holy angels also and all creatures acknowledge him, and give glory to God. For so it teaches us also in all our doings, and even in the very grievous calamities and persecutions, whereof it shall prophesy moreover, to acknowledge the providence and good will of God towards us, and his most just government. If we do this with quiet minds, we shall bear even the most heavy burdens patiently; we shall cease with curious questions to inquire why God permits Antichrist to arise, to increase, and to reign, to oppress the religion and saints of God? Then shall cease also the blasphemous mutinying of those who are not afraid to say, “God is indeed the Lord, he is almighty, he does what he will, and as he will; we are bond servants, and rather worse than bondmen. We are forced to bear whatever he will lay upon us,” \&c. As though God were unjust, and after a tyrannical fear terrible, and ruled after a carnal lust. It is most shameful to think thus, much more to speak it. This vision shall declare that God by his providence governs all things, and that his ways are just in all his ways, and his works are holy.

And first, John is prepared to receive this vision, yea, and we also are prepared in him. For when he had seen the door in heaven to be wide open, he heard withal, “Come up hither,” and so forth. It is surely a benefit beyond expression that the Lord opens heaven for us miserable, mortal men, and allows us to see what is done therein, or what he himself does there, and what his works or judgments are toward men. Let no man hereafter say that God does in heaven whatever he pleases, not passing judgment upon us who creep upon earth, and who also must suffer what we would not. For now he makes as it were an account of his works and, being assured, admits you as a witness to the matter.

And here he declares with a godly voice, what John should do, and how he should behave himself. Christ bids John ascend into heavenly places, not in body, but in mind. Therefore, our mind must be lifted up into the contemplation of heavenly things, and be purged as much as possible from earthly affections, that we may behold heavenly things with a heavenly contemplation. What shall we say, but that the example of John follows immediately? And immediately I was in the spirit: that is, in a spiritual contemplation, or ravished by the spirit into the faithful consideration of those things which were shown to me.

Now also is assembled an explanation of things that ought to be said: I will show you what things must be done hereafter. For after the scepter of God, rightly ordering or governing all things through Christ, immediately are declared the destinies of the church by seven seals and seven trumpets, within which are everywhere interwoven most comforting consolations and full of efficacy.

And first of all, before the Seals and symbolic representations are presented, there is a depiction or glimpse of God, and his most righteous judgment and governance in all things; and this is shown throughout chapters 4 and 5, so that it might prepare us for the things that will follow in chapters 6, 7, and 8. And it may seem to others, and to human judgment, to be grievous, harsh, and unjust. And the depiction or vision was as follows.

In heaven itself appeared a seat or throne of Majesty. He who sits therein holds in his right hand a book, closed with Seals. Beside him who sat was a Lamb, who takes the book and opens its Seals. And from this throne also proceeds a sevenfold Spirit, wonderfully expressing his virtues. Before the throne appears a sea of glass, bright and smooth like crystal. The throne itself rests like a wagon upon four beasts full of eyes and wings; beneath it appear all around, encircling the throne. A rainbow like an emerald goes around it. About the throne, by a circle, appear twenty-four seats, and so many elders sitting in them, crowned and in white array. This is the order of this second vision. Their roles will be explained: what the Lamb, what the beasts, what the Elders, and the other parts did. It suffices now to have touched the chiefest points of the vision, and to have given a shadowing of some of the same.

\index{Arethas of Caesarea}\index{Daniel (Book of)}\index{Ezekiel (Book of)}Secondly, we must see what each thing signifies. For a great part of the entire mystery depends on this. Regarding the manner of vision, John introduces nothing new concerning the revelation of Christ. For we read that such visions were exhibited to the prophets in a similar manner, as to Isaiah in chapter 6, to Ezekiel in chapters 1 and 11, and to Daniel in chapter 7, and so on. And a throne signifies imperial majesty and judicial administration. And because the throne is not on earth, but is seen in heaven, we shall think that God’s providence and administration of judgment are celestial, sound, most holy, and entirely free from all corruption. And upon this same throne is one sitting, sitting I say, not lying or standing. For God, the judge of all, is of a quiet mind, and is not moved by affections like men. There is no affection, injury, or unrighteousness to be thought upon in the universal government of all things. Elihu, in chapter 34 of Job, says: Far be wickedness from God, and iniquity from the Almighty. For he will render to a person according to his work, and according to the ways of each, he will reward them. For truly God will not condemn in vain, nor will the Almighty subvert judgment, and so on. And Arethas, Bishop of Caesarea, an old commentator, admonishes that the shape of a man was not attributed to him who sits in the seat, on purpose. For although afterward mention is made of a right hand holding the book, here no shape of a man is exhibited. But he says simply one sitting, he gives him no name. The reason is ready: for God, by his nature, cannot be defined, as he who is invisible and unmeasurable. According to the manner of men, human members are attributed to him, but to be understood by a trope. Moreover, when the same God appeared to the people of Israel at Sinai, they heard a voice only, but the Israelites saw no shape, as Moses witnesses in chapter 4 of Deuteronomy. Doubtless, that they should not express the incomprehensible with an image, and should commit idolatry, the great sin and wickedness. Paul in Acts 17 denies that the deity is like the forging of men, to the Romans. He ascribes to the greatest folly idols made after the shape of men, which should represent God. Of this we have spoken elsewhere. In the meantime, two precious stones are mentioned, which by their colors do, in a way, shadow the nature of our God, and admonish the godly of greater and more excellent things. A jasper is a green stone like an emerald. Green signifies the perpetuity of God, and that he quiets and keeps all things in life. But the sardine looks with a fine color like a bright red. For God dwells in inaccessible light; he is a consuming fire, and also charity itself. For the nature of stones, read Pliny, and so on.

But a rainbow environs the Throne round about; a rainbow is generally of diverse colors, but here it is of one color, and that of an emerald, namely green. The rainbow is a token of a perpetual grace and covenant made after the flood, as is declared in the 9th chapter of Genesis. And truly, the Throne of the high judge might put us wretched men in fear. Therefore, the rainbow puts us in remembrance of God’s grace, and that God who by his providence governs all things has bound himself in league to mankind, to whom, indeed, he wishes well. That league is still green, and always of force. The goodness of God towards men is perpetual. For though heaven should fall, and out of this Throne proceed most grievous thunderbolts, and calamities should fall upon us like a storm, yet is God in league with us, and loves us dearly.

Concerning the throne, twenty-four seats are seen, and twenty-four elders sit in them, as senators of the most mighty kingdom of God, and fathers of the heavenly hierarchy. This number is composed of twelve and twelve. But the twelve patriarchs signify the entire people of Israel, and the old church before Christ. The Christian church was planted and sprang up from the twelve Apostles, after the incarnation of Christ, so that this number twelve encompasses the whole church of the new people. Therefore, the entire universality of saints is assembled in heaven, triumphing with Christ their king. And therefore they are clothed in white raiment, namely, purged by Christ, and pure and clean from all corruption. They are also crowned, because they have overcome and now reign in eternal glory, truly kings and priests through Christ. The description also of their behavior admonishes that in them there is nothing lacking, but that they are truly blessed; and therefore they are shown sitting, not that they are judges or judge for Christ, but because they rest from their labors and are of most quiet and pure affections, sitting with the high judge. But what do these do? They give God no counsel as to what he should do, or by what means or way he may do this or that, but they approve his judgments. For they know all his works to be just and holy, which will immediately follow. What then shall we do? Shall it be fitting for us to inquire about the judgments of God, or prescribe what he should do or not do? I think not. You do not have in this universality of saints, all patriarchs, all judges and kings, all princes, and the whole people of God; you have among these, King Solomon himself, and the most excellent and wisest princes of the world; you have the Apostles, and men apostolic, martyrs, and the wise men of the whole universal world. Will you condemn their judgments? Following therefore their example, do not busy yourself with curious questions; praise the just judgments of God, and know that the Lord is just in all his ways, and holy in all his works. To whom be glory.
